% !TeX root = ../01_Projektarbeit/01_main.tex
% !TEX program = xelatex
\documentclass[12pt, a4paper]{report}
% Formatierungen und Packages
\usepackage{fontspec}
\usepackage[ngerman]{babel}% Language setting
\usepackage[
  a4paper,
  top=2cm,
  bottom=2cm,
  left=2.5cm,
  right=2.5cm,
  headheight=20pt,
  headsep=15pt
]{geometry} % Set page size and margins
\usepackage{graphicx} % Required for inserting images
%\usepackage{amsmath}
\usepackage{xcolor}
\usepackage{csquotes} % \enquote nutzen zum Zitieren \enquote 
\usepackage{fancyhdr}
\usepackage{setspace}
\usepackage{titlesec}
\usepackage{comment}
\usepackage{booktabs}
\usepackage{tabularx}
\usepackage{array}
\usepackage{pdfpages}
\usepackage{acronym}
\usepackage[
    backend=biber,
    style=authoryear,
    dashed=false,
    url=false
]{biblatex}
\usepackage{xurl} % sorgt für automatische Zeilenumbrüche bei zu langen Links im Quellenverzeichnis
\usepackage{tocbibind}
\usepackage{tocloft}
\usepackage{enumitem}
\usepackage{nameref}
\usepackage{chngcntr} % \counterwithout figure chapter
\usepackage{threeparttable}
\usepackage{caption}
\usepackage[colorlinks=true, allcolors=blue, hypertexnames=false]{hyperref} % allcolors=black setzen am Ende um Links zu schwärzen
% hypertexnames=false Verhindert doppelte PDF-Anker und fehlerhafte Sprungziele bei
% zurückgesetzten Seiten-, Kapitel- oder Abbildungszählern.
\usepackage{cleveref}

\DeclareDelimFormat{labelnamepunct}{\addcolon\space}
\DeclareFieldFormat{url}{\url{#1}}
\DeclareBibliographyDriver{online}{%
    \usebibmacro{bibindex}%
    \usebibmacro{begentry}%
    \printnames{author}%
    \setunit{\addcolon\space}%
    \printfield{title}%
    \setunit{\addperiod\space}%
    \printtext{In:\space}%
    \printfield{url}%
    \iffieldundef{urlyear}
    {}
    {~\mkbibparens{%
     \printfield{urlday}.%
     \printfield{urlmonth}.%
     \printfield{urlyear}%
   }}%
    \usebibmacro{finentry}%
}
%\setmainfont{Arial}
\setlength{\headheight}{17.01598pt} % Fehler hat gesagt Headheight größer
\addtolength{\topmargin}{-2.01598pt} % Befehl ist um Text weiter hoch zu schieben wegen \headheight{17.01598pt} (es müsste eigentlich 15pt sein
\setstretch{1.3}
\titleformat{\chapter}[hang]
    {\fontsize{14pt}{22pt}\selectfont\bfseries} % Format
    {\thechapter.} % Label
    {1em} %Abstand
    {} % <Vor-Code>
\titleformat{\section}[hang]
    {\fontsize{14pt}{18pt}\bfseries} % Format
    {\thesection} % Label
    {1em} %Abstand
    {} % <Vor-Code>
\titleformat{\subsection}[hang]
    {\fontsize{12pt}{18pt}\bfseries} % Format
    {\thesubsection} % Label
    {1em} %Abstand
    {} % <Vor-Code>
\titlespacing*{\chapter}{0pt}{-26pt}{12pt}
\DeclareNameFormat{nurNachname}
  {\MakeUppercase{\namepartfamily}}
  
\DeclareCiteCommand{\footdirekt}[\mkbibfootnote] % Zitat direkt
  {}
  { \printnames[nurNachname]{author},
    \printfield{year},~\printfield{postnote}\nopunct}
  {}
  {}
\DeclareCiteCommand{\footvgl}[\mkbibfootnote] % Vergleich Zitat
  {}
  { vgl. \printnames[nurNachname]{author}, \printfield{year},~\printfield{postnote}\nopunct}
  {}
  {}
\DeclareCiteCommand{\footonline}[\mkbibfootnote]
  {}
  {
    online:
    \printnames{author},
    \printdate,
    \mkbibparens{%
      \printfield{urlday}.%
      \printfield{urlmonth}.%
      \printfield{urlyear}%
    }%
    \nopunct
  }
  {}
  {}
\newcommand{\abschnitt}[1]{%
  \paragraph{#1}\mbox{}\par
}
\counterwithout{figure}{chapter} % Kapitelnummer bei Abbildungen entfernen, fortlaufende Nummerierung
\counterwithout{table}{chapter} % Kapitelnummer bei Tabellen entfernen, fortlaufende Nummerierung
%\captionsetup[figure]{ Ausgeklammert da sonst ein Fehler kommt wenn man es nicht verwendet
%    labelfont=bf,
%    labelsep=space,
%   justification=raggedright,
%    singlelinecheck=false,
%    format=hang
%}
\addbibresource{literatur.bib}
\begin{document}
\setlength{\parindent}{0pt} % Einzug nach Bildern im gesamten Dokument deaktivieren
\fancypagestyle{projektarbeit}{
    \fancyhf{} % Kopf- und Fußzeile leeren
    \fancyhead[R]{Paul Lorenz}
    \fancyhead[L]{\nouppercase{Kopie der Textinhalte aus Internetquellen}}
    \fancyfoot[R]{\thepage}
}
\fancypagestyle{plain}{
    \fancyhf{} % Kopf- und Fußzeile leeren
    \fancyhead[R]{Paul Lorenz}
    \fancyhead[L]{\nouppercase{\leftmark}}
    \fancyfoot[R]{Seite \thepage}
}
\renewcommand{\chaptermark}[1]{% Kommentar bis zum Zeilenende“. Hier werden sie verwendet, um unerwünschte Leerzeichen zu verhindern.
  \markboth{\thechapter.\ #1}{}%
} % Entfernen des Wort Kapitels in der Kopfzeile
\pagestyle{projektarbeit}
\pagenumbering{arabic} % Arabische Numerierung der Seitenzahlen
\abschnitt{Kopie der Textinhalte aus Internetquellen}
\vspace{\baselineskip}
\fullcite{SAP-Geschichte}

{\small
\enquote{Am 1. April 1972 gründen fünf ehemalige IBM-Mitarbeiter das Unternehmen „Systemanalyse Programmentwicklung“. Die Vision von Dietmar Hopp, Hasso Plattner, Claus Wellenreuther, Klaus Tschira und Hans-Werner Hector: Standardsoftware für Unternehmen zu entwickeln, die alle betrieblichen Abläufe integriert und es ermöglicht, Daten in Echtzeit zu verarbeiten.}
\vspace{\baselineskip} % Das erzeugt eine Leerzeile
}

\fullcite{SAP-Erklaerung}

{\small
\enquote{Ein ERP-System besteht aus integrierten Modullösungen oder Geschäftsanwendungen, die sich eine gemeinsame Datenbank teilen, durch die sie miteinander verbunden sind und miteinander kommunizieren können. Jedes ERP-Modul ist üblicherweise auf einen Geschäftsbereich ausgerichtet. Sie greifen aber ineinander und nutzen dieselben Daten, um die Anforderungen des Unternehmens zu erfüllen.

Unternehmen wählen die gewünschte Modullösung aus – z. B. Finanzen, Logistik, Beschaffung und Personalwesen – und können sie je nach Bedarf ergänzen und skalieren. ERP-Systeme können auch branchenspezifische Anforderungen abdecken, entweder als Teil der Kernfunktionen des Systems oder durch Anwendungserweiterungen, die sich in die Modulsuite einfügen.}
\vspace{\baselineskip} % Das erzeugt eine Leerzeile
}

\fullcite{SAP-PCo}

{\small
\enquote{Mit PCo können Sie automatisch oder auf Anforderung Datenpunkte (Tags) und Ereignisse aus den angeschlossenen Quellsystemen in der Produktion empfangen und diese an die angeschlossenen SAP-Systeme weitergeben. Weiterhin können Sie über PCo OPC UA-Methodenaufrufe absetzen, entgegennehmen und verarbeiten und somit auch komplexe Koordinierungsaufgaben im Umfeld von Maschinensteuerungen (SPS) abbilden. Der PCo Web-Server erweitert diese Möglichkeiten und erlaubt die Bereitstellung von selbst-definierten Methoden, die als Web-Service-Operationen aufgerufen werden können.}
\vspace{\baselineskip} % Das erzeugt eine Leerzeile
}

\fullcite{S/4HANA}

{\small
\enquote{SAP S/4HANA ist eine Enterprise-Resource-Planning-Lösung, die auf einer \linebreak
SAP-HANA-Datenbank basiert und einen On-Premise-Betrieb ermöglicht. Das bedeutet, dass Sie als Kunde sich um die Installation, das Upgrade und den Betrieb der Software in der Infrastruktur Ihres Unternehmens kümmern.}
}

\fullcite{Kommentar}

{\small
\enquote{Die Anwendung des neuen Algorithmus wird uns dabei helfen den MES-Implementierungsprozess weiter zu standardisieren. Daraus resultiert eine schnelle und effiziente Anbindung neuer Anlagen und somit eine Optimierung des gesamten Rollout-Prozesses.}}


\end{document}



